\documentclass[11pt,a4paper]{article}

\usepackage[margin=1in]{geometry}
\usepackage{amsmath,amssymb,amsthm}
\usepackage{enumitem}
\usepackage{booktabs}
\usepackage{hyperref}
\usepackage{xcolor}
\usepackage{tikz}
\usepackage{algorithm}
\usepackage{algpseudocode}
\usepackage{microtype}

\hypersetup{
  colorlinks=true,
  linkcolor=blue!50!black,
  urlcolor=blue!60!black,
  citecolor=blue!50!black,
}

\newtheorem{theorem}{Theorem}
\newtheorem{lemma}[theorem]{Lemma}
\newtheorem{proposition}[theorem]{Proposition}
\newtheorem{corollary}[theorem]{Corollary}
\theoremstyle{definition}
\newtheorem{definition}{Definition}
\theoremstyle{remark}
\newtheorem{remark}{Remark}

\newcommand{\name}{\mathsf{name}}
\newcommand{\rec}{\mathsf{rec}}
\newcommand{\KEM}{\mathsf{KEM}}
\newcommand{\Sig}{\mathsf{Sig}}
\newcommand{\Adv}{\mathsf{Adv}}
\newcommand{\negl}{\mathsf{negl}}
\newcommand{\KGen}{\mathsf{KGen}}
\newcommand{\Encap}{\mathsf{Encap}}
\newcommand{\Decap}{\mathsf{Decap}}
\newcommand{\Sign}{\mathsf{Sign}}
\newcommand{\Verify}{\mathsf{Verify}}
\newcommand{\pk}{\mathit{pk}}
\newcommand{\sk}{\mathit{sk}}
\newcommand{\RNS}{\textsc{ZAP-RNS}}

\title{\RNS: Cryptographic Name--Identity Binding for In-Cluster Service Discovery}
\author{ZAP Protocol Authors\\\texttt{zap-proto.io}}
\date{2026}

\begin{document}
\maketitle

\begin{abstract}
We present \RNS, a service-naming scheme in which every name is bound at
registration time to a long-term key encapsulation mechanism (KEM) public
key. Claims about a name --- its existence, its current public key, its
metadata --- are signed under the corresponding private key, and lookups
return a verifiable chain rooted at a per-cluster authority public key.
This eliminates name spoofing and Sybil attacks at the protocol layer:
spoofing the name $N$ reduces to recovering the KEM private key for $N$
or forging a signature under it. We compare \RNS{} against DNS, DNSSEC,
mDNS, and SPIFFE/SPIRE; give a formal model in which the
name-spoofing advantage is bounded by KEM key-recovery and signature
forgery probabilities; and discuss the operational properties of
caching, revocation, and federation. \RNS{} is post-quantum native: its
binding primitive is a KEM (ML-KEM-768 by default) and its signature
primitive is composable with ML-DSA-65, so the scheme inherits PQ
guarantees without retrofit. We argue \RNS{} is the right primitive for
in-cluster service mesh discovery because identity is checked at every
lookup rather than only at TLS handshake, and because the trust root is
small, audit-able, and pluggable into existing workload-attestation
infrastructure.
\end{abstract}

\section{Introduction}

Service discovery has accumulated layers of compromise. The Domain Name
System~\cite{rfc1035} answers ``where is $N$?'' but not ``is the holder
of $N$ who I think it is?'' DNSSEC~\cite{rfc4033} signs records but
deploys end-to-end in well under 30\% of the public Internet and almost
never inside a cluster. Multicast DNS~\cite{rfc6762} and SRV
records~\cite{rfc2782} push naming further into the local network with
no identity binding at all: any host on the segment can claim any name.

Cluster-scale systems --- ZooKeeper, etcd, Consul, Kubernetes ---
collapse name and address into a key/value store guarded by an out-of-band
authentication tier. Discovery is done by querying the store; the
identity of the discovered service is not part of the record but
re-established later via TLS, mTLS, or workload attestation. This is
the SPIFFE/SPIRE model~\cite{spiffe}: a name (URI SVID) is issued as
an X.509 certificate by a workload attestor on the assumption that the
Kubernetes API server, kubelet, and node identity are trusted. The
identity binding lives in a separate plane from the discovery plane.

The recurring failure across these designs is the same: \emph{the name
itself does not commit to a key}. To go from name to authenticated
endpoint, every system relies on at least one out-of-band step ---
DNSSEC chain, X.509 chain, mesh CA, SPIFFE attestor --- and the
operational burden of those steps is what causes them to be skipped
in practice.

We propose \RNS, a scheme in which the name \emph{is} a commitment to a
public key. Registration produces a record
$(N, \pk_N^\KEM, \pk_N^\Sig, \mathit{ttl}, \sigma)$ where $\sigma$ is a
chain of signatures terminating at a single per-cluster RNS-root
public key. Lookups return the full record. Connections encrypt to
$\pk_N^\KEM$, so reaching an impostor requires either knowing
$\sk_N^\KEM$ or compromising the RNS root. The same record is
verifiable offline by any client that has been provisioned with the
root once. Identity, naming, and key distribution become a single
primitive.

The contributions of this paper are:
\begin{enumerate}[label=(\arabic*),leftmargin=*]
\item A formal definition of \RNS{} records, registration, and lookup
(\S\ref{sec:def}).
\item A security model and reductions showing that name-spoofing
advantage is bounded by KEM key-recovery, signature forgery, and
RNS-root compromise (\S\ref{sec:sec}, \S\ref{sec:formal}).
\item A Sybil-resistance argument tying registration cost to the
cluster's chosen workload-attestation primitive (\S\ref{sec:sybil}).
\item A comparison of \RNS{} with DNS, DNSSEC, mDNS, and SPIFFE on
six axes including post-quantum readiness (\S\ref{sec:cmp}).
\item Implementation notes on caching, revocation, and multi-authority
federation (\S\ref{sec:impl}).
\end{enumerate}

\section{Definition of \RNS}\label{sec:def}

Let $\KEM = (\KEM.\KGen, \KEM.\Encap, \KEM.\Decap)$ be an IND-CCA2 secure
key encapsulation mechanism and let $\Sig = (\Sig.\KGen, \Sig.\Sign,
\Sig.\Verify)$ be an EUF-CMA secure signature scheme. Default choices
are ML-KEM-768~\cite{mlkem} and ML-DSA-65~\cite{mldsa}, but the scheme
is parametric in the primitives.

\begin{definition}[Record]\label{def:record}
A \RNS{} record is a tuple
\[
\rec = \bigl(N,\ \pk^\KEM,\ \pk^\Sig,\ t_0,\ \mathit{ttl},\ m,\ \sigma_{\mathrm{owner}},\ \sigma_{\mathrm{auth}}\bigr)
\]
where $N \in \{0,1\}^*$ is the human-readable name, $\pk^\KEM$ and
$\pk^\Sig$ are the long-term KEM and signing public keys, $t_0$ is
issuance time, $\mathit{ttl}$ is validity duration, $m$ is auxiliary
metadata (endpoint, attestation evidence), $\sigma_{\mathrm{owner}} =
\Sig.\Sign(\sk^\Sig, H(N \| \pk^\KEM \| \pk^\Sig \| t_0 \| \mathit{ttl}
\| m))$ is the self-signed binding, and $\sigma_{\mathrm{auth}} =
\Sig.\Sign(\sk_{\mathrm{root}}, H(\rec \setminus \sigma_{\mathrm{auth}}))$
is the authority counter-signature.
\end{definition}

\paragraph{Trust root.} A cluster operator generates a single
$(\pk_{\mathrm{root}}, \sk_{\mathrm{root}})$ pair. $\pk_{\mathrm{root}}$
is bundled into every workload's image or sidecar at provisioning time.
$\sk_{\mathrm{root}}$ lives in an HSM or threshold service. There is
exactly one root per trust domain, and it is the single point of
human-audit-able trust.

\paragraph{Registration.} A workload runs the protocol of
Algorithm~\ref{alg:register}. It generates KEM and signing keypairs,
constructs an unsigned record, signs it under $\sk^\Sig$, and submits
the record together with attestation evidence (e.g., a Kubernetes
ServiceAccount token, a TPM quote, a TEE attestation) to the RNS
authority. The authority validates evidence under its policy, checks
that $N$ is not already registered (or that the requester is authorised
to roll the key), and emits $\sigma_{\mathrm{auth}}$.

\begin{algorithm}
\caption{Registration}\label{alg:register}
\begin{algorithmic}[1]
\State $(\pk^\KEM, \sk^\KEM) \gets \KEM.\KGen()$
\State $(\pk^\Sig, \sk^\Sig) \gets \Sig.\KGen()$
\State $r \gets (N, \pk^\KEM, \pk^\Sig, t_0, \mathit{ttl}, m)$
\State $\sigma_{\mathrm{owner}} \gets \Sig.\Sign(\sk^\Sig, H(r))$
\State submit $(r, \sigma_{\mathrm{owner}}, \mathit{evidence})$ to RNS authority
\State authority verifies evidence and $\Sig.\Verify(\pk^\Sig, H(r), \sigma_{\mathrm{owner}})$
\State authority emits $\sigma_{\mathrm{auth}}$ and stores
$\rec = (r, \sigma_{\mathrm{owner}}, \sigma_{\mathrm{auth}})$
\end{algorithmic}
\end{algorithm}

\paragraph{Lookup.} A client queries the RNS authority for $N$ and
receives $\rec$. Verification:
\begin{enumerate}[leftmargin=*]
\item Check $t_0 \le \mathrm{now} \le t_0 + \mathit{ttl}$.
\item Check $\Sig.\Verify(\pk^\Sig, H(r), \sigma_{\mathrm{owner}})$.
\item Check $\Sig.\Verify(\pk_{\mathrm{root}}, H(r \| \sigma_{\mathrm{owner}}), \sigma_{\mathrm{auth}})$.
\item Consult the freshness oracle (revocation list or short-lived
witness) for $\rec$.
\end{enumerate}
On success the client treats $\pk^\KEM$ as the only legitimate
encapsulation target for $N$.

\paragraph{Connection.} The client encapsulates a fresh shared secret
under $\pk^\KEM$ and sends the ciphertext together with its first
record. Only the holder of $\sk^\KEM$ for $N$ can decapsulate.
Subsequent application traffic uses the derived secret in an AEAD.
Identity is checked \emph{on every lookup}, not only at handshake.

\paragraph{Encrypted lookup.} The lookup itself can be encapsulated
to $\pk_{\mathrm{root}}$ (or a sub-key derived from it via HPKE), so
that passive observers see only an encrypted blob between client and
RNS. This degrades worst-case adversary capability from ``learns who
talks to whom'' to ``learns who talks to RNS''.

\section{Security Properties}\label{sec:sec}

\paragraph{Authenticity.} A name not registered in RNS cannot be
served. A client receiving a record for $N$ verifies a signature chain
to $\pk_{\mathrm{root}}$; absent the root key, no adversary can
fabricate a record that verifies.

\paragraph{Continuity.} Once $N$ is bound to $\pk_N^\KEM$, the binding
changes only via a fresh authority counter-signature. The authority
policy decides whether re-binding requires (a) a signature from the
old $\sk^\Sig$ (Lamport-style continuity~\cite{lamport}), (b) re-running
the original attestation, or (c) explicit operator approval. \RNS{}
encodes the policy choice in $m$ and refuses lookups whose continuity
chain breaks the policy.

\paragraph{Privacy.} With encrypted lookup, a network observer learns
neither the names a client resolves nor the public keys returned. The
authority still observes both, mitigated by classic
PIR~\cite{chor1995pir} or OHTTP-style~\cite{ohttp} interposition. We
treat scalable privacy-preserving lookup as an open problem
(\S\ref{sec:open}).

\paragraph{Freshness and revocation.} TTLs are short (default
$300\,\mathrm{s}$). Each authority publishes a signed Merkle witness
of currently-valid records every TTL/2; clients consult it at each
lookup. Compromised KEM keys are revoked by appending a tombstone
record signed by the authority.

\section{Formal Model}\label{sec:formal}

We define the name-spoofing experiment $\mathbf{Exp}^{\mathrm{spoof}}_{\RNS}(\mathcal{A})$:

\begin{enumerate}[leftmargin=*,label=\arabic*.]
\item Challenger generates $(\pk_{\mathrm{root}}, \sk_{\mathrm{root}})$
and gives $\pk_{\mathrm{root}}$ to $\mathcal{A}$.
\item $\mathcal{A}$ may query a registration oracle on names of its
choice; the oracle runs Algorithm~\ref{alg:register} honestly and
returns the resulting records and the public keys, but \emph{not} the
private keys.
\item $\mathcal{A}$ outputs a target name $N^*$ (not previously
registered by $\mathcal{A}$ as the legitimate owner) and a forged
record $\rec^*$.
\item $\mathcal{A}$ wins if $\rec^*$ verifies against $\pk_{\mathrm{root}}$
and $\rec^*$'s public key differs from the public key on file at the
challenger for $N^*$ (or no record exists for $N^*$).
\end{enumerate}

Theorems~1--3 of the companion proof file bound
\[
\Adv^{\mathrm{spoof}}_{\RNS}(\mathcal{A}) \;\le\;
\Adv^{\mathrm{EUF\text{-}CMA}}_{\Sig}(\mathcal{B}_1) \;+\;
\Adv^{\mathrm{key\text{-}rec}}_{\KEM}(\mathcal{B}_2) \;+\;
\Pr[\text{root compromise}],
\]
where the third term is policy-controlled and audit-able.

\paragraph{Connection binding.} Theorem~2 of the proof file shows
that under IND-CCA2 of $\KEM$, a client encrypting to name $N^*$
reaches a peer holding $\sk_{N^*}^\KEM$ with overwhelming probability,
or learns nothing.

\section{Sybil Resistance}\label{sec:sybil}

\RNS{} reduces Sybil resistance to attestation-bypass cost. The
authority emits $\sigma_{\mathrm{auth}}$ only after attestation evidence
satisfies its policy. The policy is pluggable:
\begin{itemize}[leftmargin=*]
\item \textbf{K8s ServiceAccount.} The cluster's API server is the
attestor; cost = stealing a ServiceAccount token.
\item \textbf{TPM quote.} A measured-boot TPM signs a PCR digest;
cost = compromising a hardware root of trust.
\item \textbf{TEE attestation.} An SEV-SNP / TDX / SGX quote;
cost = breaking the TEE.
\item \textbf{Proof-of-stake.} Each registration burns or stakes a
token; cost = direct economic.
\end{itemize}

The Sybil bound is tight: an adversary capable of $n$ legitimate
registrations under the policy can bind $n$ names. The cluster operator
chooses the per-registration cost to match the consequences. There is
no ambient Sybil channel because the only way to insert a record is
through the attestor.

\section{Comparison}\label{sec:cmp}

Table~\ref{tab:cmp} compares discovery schemes on six properties:
intrinsic name-to-key binding, lookup-time authentication of records,
encryption-to-name (a client can encrypt directly to a returned
identity), post-quantum readiness, deployment effort relative to a bare
DNS deployment, and resistance to a passive network observer learning
the names being resolved.

\begin{table}[h]
\centering
\small
\begin{tabular}{lcccccc}
\toprule
& \textbf{Name binds key} & \textbf{Auth lookup} & \textbf{Enc-to-name} & \textbf{PQ ready} & \textbf{Op cost} & \textbf{Lookup priv.} \\
\midrule
DNS~\cite{rfc1035}      & no    & no    & no    & n/a & low    & no \\
DNSSEC~\cite{rfc4033}   & no\textsuperscript{*} & yes & no & retro & high & no \\
mDNS~\cite{rfc6762}     & no    & no    & no    & n/a & low    & no \\
SRV~\cite{rfc2782}      & no    & no    & no    & n/a & low    & no \\
SPIFFE~\cite{spiffe}    & X.509 & at TLS only & no & retro & high & no \\
HIP~\cite{rfc7401}      & yes (HI) & yes & yes & retro & high & partial \\
Tor onion~\cite{tor}    & yes (v3) & yes & yes & no & high & yes \\
\textbf{\RNS} (this work) & \textbf{yes} & \textbf{yes (every lookup)} & \textbf{yes} & \textbf{native} & \textbf{med} & \textbf{yes (encrypted)} \\
\bottomrule
\end{tabular}
\caption{Comparison of service discovery schemes. ``Name binds key'':
the name is a commitment to a public key. ``Auth lookup'': lookups
return signed records that any client can verify. ``Enc-to-name'':
clients can encrypt directly to the returned identity without a
preceding handshake. \textsuperscript{*}DNSSEC binds the \emph{zone}
to a key, not individual service names to service keys.}
\label{tab:cmp}
\end{table}

\RNS{} is the only scheme that ticks every column. The closest
analogue is the Host Identity Protocol (HIP)~\cite{rfc7401}, which
also separates a host identity from its locator and binds names to
public keys. \RNS{} differs in four ways: (i) it is post-quantum native;
(ii) lookups return a chain to a single auditable root, not arbitrary
DNS resource records; (iii) registration is gated by pluggable
attestation rather than the open Internet PKI; (iv) the binding is to
a KEM key, enabling encryption-to-name without a separate key-exchange
roundtrip.

\section{Implementation Notes}\label{sec:impl}

\paragraph{Caching.} Records are immutable for their TTL. Clients cache
verified records and a freshness witness; a cache hit costs one
signature verification. Negative caching of NXDOMAIN responses, signed
by the authority, prevents downgrade-to-no-record attacks.

\paragraph{Revocation.} Short TTLs ($60\,\mathrm{s}$ to
$300\,\mathrm{s}$) replace explicit revocation in the common case. For
emergency revocation the authority publishes a signed tombstone with
the same name, which evicts cached records on next freshness check.

\paragraph{Federation.} Multiple authorities co-exist by being placed
in a forest: each authority $A_i$ has $\pk_{\mathrm{root},i}$, and
clients hold a known-roots set. Names are namespaced by authority
($N@A_i$), so collisions are impossible. Cross-authority delegation
is a counter-signature on a delegation record, identical in shape to
$\sigma_{\mathrm{auth}}$.

\paragraph{Operational primitives.} The authority is a small stateful
service: $\sim$10\,kB per record, $O(1)$ verification per lookup, and
no external dependencies beyond the HSM holding $\sk_{\mathrm{root}}$.
For a cluster of $10^5$ services the on-disk state is below 1\,GB and
fits in memory on a single node.

\paragraph{Sidecar integration.} A workload's sidecar exposes a local
UDS, performs registration on startup, caches its own KEM private key
in a tmpfs (or HSM), and intercepts outgoing connections to translate
$N$ into $(\pk_N^\KEM, \mathrm{endpoint})$. No application code change
is required; legacy clients see a familiar resolver interface.

\section{Open Questions}\label{sec:open}

\paragraph{Lookup privacy at scale.} Encrypting the lookup to the
authority hides names from network observers, but the authority itself
still sees them. For large public clusters one wants a private
information retrieval~\cite{chor1995pir} or oblivious-HTTP~\cite{ohttp}
overlay. Both are quadratic-or-worse in the worst case. We have not
benchmarked PIR-over-RNS in this paper.

\paragraph{Revocation latency.} Short TTLs reduce revocation latency
but inflate authority load. We conjecture an adaptive TTL keyed to a
service's request rate is workable, but have not analysed the
adversarial-load model.

\paragraph{Multi-region authority placement.} Federated authorities
trade global consistency for partition tolerance. CRDT-based
authority sets are an active area; we leave them future work.

\paragraph{Hybrid signatures.} ML-DSA-65 is a sound default, but a
hybrid (ML-DSA + Ed25519) registration signature would tolerate a
catastrophic break in either scheme. The cost is roughly doubled
record size; we treat this as a configuration parameter.

\section{Conclusion}

\RNS{} makes the name into a key. The reduction is sharp: spoofing
$N$ requires breaking a KEM, forging a signature, or compromising a
single audit-able root. The scheme is post-quantum native, integrates
with any workload-attestation primitive a cluster already runs, and
adds at most one signature verification per lookup beyond a
DNS-style cache. For in-cluster service mesh discovery --- where the
trust domain is small, attestation is already required, and identity
must be checked at every hop --- \RNS{} is a strict improvement on the
status quo of layered DNS, X.509, and SPIFFE.

\bibliographystyle{plain}
\bibliography{refs}

\end{document}
